\documentclass{beamer}
\usepackage{amsmath}
\setbeameroption{show notes on second screen=right}
\setbeamertemplate{note page}[plain]

\begin{document}

\begin{frame}
\frametitle{}

When adding or subtracting fractions, the two fractions that are involved need a common denominator. 
\note[item]{\alert<1>{When adding or subtracting fractions, the two fractions that are involved need a common denominator.}}
\pause (When multiplying or dividing fractions, there is no requirement for a common denominator.) 
\note[item]{\alert<2>{When multiplying or dividing fractions, there is no requirement for a common denominator.}}

\vskip10pt

\pause 
{\bf Example.} $\frac23+\frac37$. \note[item]{\alert<3>{Let's look at an example. What is two thirds plus three sevenths?}}

\pause\note[item]{\alert<4>{We will rewrite both fractions to have twenty-one in the denominator.}}

\vskip4pt

\pause $=\frac{2 \cdot 7}{3 \cdot 7} + \pause \frac{3 \cdot 3}{7 \cdot 3}$ 
\note[item]{\alert<5>{For the first fraction, multiply the top and the bottom by seven. We haven't finshed this step, which explains the space after the plus sign.}}

\vskip4pt

\note[item]{\alert<6>{For the second fraction, multiply the top and the bottom by three.}}
\pause $=\frac{14}{21} + \frac{9}{21}$
\note[item]{\alert<7>{Simplifying each numerator and denominator separately, we have fourteen over twenty-one plus nine over twenty-one.}}

\vskip4pt

\pause $=\frac{23}{21}$
\note[item]{\alert<8>{Now that we have the same denominator, we only add numerators. Since fourteen plus nine is twenty-three, our single fraction is twenty-three over twenty-one. We prefer improper fractions over mixed fractions.}}

\end{frame}

\begin{frame}
\frametitle{}

{\bf Example.} $\frac{1}{10}+\frac{1}{6}$. 
\note[item]{\alert<1>{Let's look at another example. What is one tenth plus one sixth?}}

\vskip10pt

\pause 
\note[item]{\alert<2>{We could rewrite both fractions to have sixty in the denominator, but ten and six both have a common factor of two, so will rewrite both fractions to have thirty in the denominator.}}

\vskip4pt

\pause $=\frac{1 \cdot 3}{10 \cdot 3} + \pause \frac{1 \cdot 5}{6 \cdot 5}$
\note[item]{\alert<3>{For the first fraction, multiply the top and the bottom by three.}}
\note[item]{\alert<4>{For the second fraction, multiply the top and the bottom by five.}}

\vskip4pt

\pause $=\frac{3}{30} + \frac{5}{30}$
\note[item]{\alert<5>{Simplifying each numerator and denominator separately, we have three over thirty plus five over thirty.}}

\vskip4pt

\pause $=\frac{8}{30}$
\note[item]{\alert<6>{Now that we have the same denominator, we only add numerators. Since three plus five is eight, our single fraction is eight over thirty.}}

\vskip4pt

\pause $=\frac{4}{15}$
\note[item]{\alert<7>{We can simplify this fraction by cancelling a factor of two on top and bottom, leaving us with four over fifteen.}}

\end{frame}


\begin{frame}
\frametitle{}

{\bf Example.} $\frac{3}{14}+\frac{1}{7}$. 
\note[item]{\alert<1>{Let's look at another example. What is three fourteenths plus one seventh?}}

\vskip10pt

\pause
\note[item]{\alert<2>{Since fourteen times seven is ninety-eight, we could rewrite both fractions to have ninety-eight in the denominator, but we can actually just make both fractions have fourteen in the denominator.}}

\pause $=\frac{3}{14} + \pause \frac{1 \cdot 2}{7 \cdot 2}$ \pause
\note[item]{\alert<3>{Copy the first fraction as is, with no changes.}}
\note[item]{\alert<4>{For the second fraction, multiply the top and the bottom by two.}}
\note[item]{\alert<5>{Unlike in previous examples, we didn't even change how both fractions are written. We only rewrote the second fraction.}}

\vskip4pt
\pause $=\frac{3}{14} + \frac{2}{14}$
\note[item]{\alert<6>{Simplifying the second numerator and denominator separately, we have three over fourteen plus two over fourteen.}}

\vskip4pt
\pause $=\frac{5}{14}$
\note[item]{\alert<7>{Now that we have the same denominator, we only add numerators. Since three plus two is five, our single fraction is five over fourteen.}}

\end{frame}


%% Blank frames at the end to prevent weird shifts.
\begin{frame}
\end{frame}
\begin{frame}
\end{frame}

\end{document}
